\documentclass[../main.tex]{subfiles}


\begin{document}

\ifSubfilesClassLoaded{

    \tableofcontents
    
}{}

\chapter{Measure Theory}

\section{Outer Measure}
\begin{definition}{Outer Measure}{}
A function $\varphi$ defined for every subset $A$ of an arbitrary set $X$ is called an \dfntxt{outer measure} on $X$ if the following conditions are satisfied:
\begin{enumerate}
    \item[(i)] $\varphi(\emptyset) = 0$,
    \item[(ii)] $0 \leq \varphi(A) \leq \infty$ whenever $A \subset X$,
    \item[(iii)] $\varphi(A_1) \leq \varphi(A_2)$ whenever $A_1 \subset A_2$,
    \item[(iv)] $\varphi(\bigcup_{i=1}^{\infty} A_i) \leq \sum_{i=1}^{\infty} \varphi(A_i)$ for every countable collection of sets $\{A_i\}$ in $X$.
\end{enumerate}
\end{definition}

\begin{example}{Outer Measure}{}
\begin{enumerate}
\item[(i)] In an arbitrary set $X$, define $\varphi(A) = 1$ if $A$ is nonempty and $\varphi(\emptyset) = 0$.
\item[(ii)] Let $\varphi(A)$ be the number (possibly infinite) of points in $A$.
\item[(iii)] Let $\varphi(A)=\left\{\begin{array}{ll} 0 & \text { if card } A \leq \aleph_{0}, \\ 1 & \text { if card } A > \aleph_{0} . \end{array}\right.$
\item[(iv)] If $X$ is a metric space, fix $\varepsilon > 0$. Let $\varphi(A)$ be the smallest number of balls of radius $\varepsilon$ that cover $A$.
\item[(v)] Select a fixed $x_0$ in an arbitrary set $X$, and let
\[
\varphi(A)=\left\{\begin{array}{ll}
0 & \text { if } x_{0} \notin A, \\
1 & \text { if } x_{0} \in A;
\end{array}\right.
\]
$\varphi$ is called the \dfntxt{Dirac measure} concentrated at $x_0$.
\end{enumerate}
\end{example}

\begin{definition}{}{}
Let $\varphi$ be an outer measure on a set $X$. A set $E \subset X$ is called \dfntxt{$\varphi$-measurable} if
\[
\varphi(A) = \varphi(A \cap E) + \varphi(A - E)
\]
for every set $A \subset X$. In view of property (iv) above, observe that $\varphi$-measurability requires only
\[
\varphi(A) \ge \varphi(A \cap E) + \varphi(A - E).
\]
\end{definition}

\begin{note}
    也就是说, \(  E  \)是可测的, 当且仅当任意集合 \(  A  \)都可以被 \(  E  \)按外测度被良好地拆解为 \(  A\cap E  \)和 \(  A\setminus E  \).      
\end{note}

\begin{lemma}{}{7-23-1}
A set $E \subset X$ is $\varphi$-measurable if and only if
\[
\varphi(P \cup Q) = \varphi(P) + \varphi(Q)
\]
for all sets $P$ and $Q$ such that $P \subset E$ and $Q \subset E^{c}$.
\end{lemma}

\begin{note}
    也就是说 \(  E  \)可测, 当且仅当 落在\(  E  \)的内, 外的集合 在外测度意义下是无关的. 
\end{note}

\begin{proof}
   For Necessity, if \(  E   \) is \(   \varphi   \)-measurable , then for all such \(  P,Q  \),  \[
     \varphi \left( P\cup Q \right)=  \varphi \left( \left( P\cup Q \right)\cap E  \right) +  \varphi \left( \left( P\cup Q \right)-E  \right)=  \varphi \left( P \right)+  \varphi \left( Q \right)     
    \]  

    For sufficiency ,  let \(  P= A\cap E  \) , \(  Q= A\cap E^{c}= A\setminus  E  \) .     Then \[
       \varphi \left( A \right)=  \varphi \left( P\cup Q \right)=   \varphi \left( P \right)+  \varphi \left( Q \right)    =\varphi \left( A\cap E \right)+  \varphi \left( A\setminus E \right)  
    \]
    
    
\end{proof}

\begin{theorem}{}{7-24-1}
Suppose $\varphi$ is an outer measure on an arbitrary set $X$.
Then the following four statements hold:
\begin{enumerate}
\item[(i)] $E$ is $\varphi$-measurable whenever $\varphi(E) = 0$.
\item[(ii)] $\emptyset$ and $X$ are $\varphi$-measurable.
\item[(iii)] $E_1 - E_2$ is $\varphi$-measurable whenever $E_1$ and $E_2$ are $\varphi$-measurable.
\item[(iv)] If $\{E_i\}$ is a \dfntxt{countable collection of disjoint $\varphi$-measurable sets}, then $\bigcup_{i=1}^\infty E_i$ is $\varphi$-measurable and
\[
\varphi\left(\bigcup_{i=1}^\infty E_i\right) = \sum_{i=1}^\infty \varphi(E_i).
\]
\end{enumerate}
More generally, if $A \subset X$ is an arbitrary set, then
\[
\varphi(A) = \sum_{i=1}^\infty \varphi(A \cap E_i) + \varphi(A \cap S^{c}),
\]
where
\[
S = \bigcup_{i=1}^\infty E_i.
\]
\end{theorem}
\begin{proof}
    \begin{enumerate}
        \item Take any \(  A\subseteq X  \) . Since \(  A\cap E\subseteq E  \) , \(   \varphi \left( A\cap E \right)= 0   \) .  Thus \[
         \varphi \left( A\cap E \right)+   \varphi \left( A\cap E^{c} \right)=  \varphi \left( A\cap E^{c} \right)\le  \varphi \left( A \right)    
        \]   
        \item Follows immediately frome lemma \ref{lem:7-23-1}.
        \item Take any \(  P,Q  \) such that \(  P\subseteq \left( E_1\setminus E_2 \right) = E_1\cap  E_2^{c}  \) , \(  Q\subseteq \left( E_1\setminus E_2 \right)^{c}= E_1^{c}\cup E_2   \) .   \[
        \begin{aligned}
         \varphi \left( Q \right)&=  \varphi \left( Q\cap E_2 \right)+  \varphi \left( Q-  E_2 \right)
        \end{aligned}   
        \] Note that \(  Q-E_2 \subseteq E_1^{c}  \) , \(  Q\cap E_2\subseteq E_2  \) . Since \(  P\subseteq E_1  \) , the measurability of \(  E_1  \) implis that  \[
         \varphi \left( P \right)+  \varphi \left( Q-E_2 \right)=  \varphi \left( P\cup \left( Q-E_2 \right)  \right)   
        \]  Note that \(  P\cup \left( Q-E_2 \right)\subseteq E_2^{c}   \)  . Since \(  Q\cap E_2\subseteq E_2  \) , we have \[
         \varphi \left( P\cup \left( Q-E_2 \right) \right)+  \varphi \left( Q\cap E_2 \right)=  \varphi \left( P\cup Q \right)   
        \]  Finally , we have \[
        \begin{aligned}
         \varphi \left( P \right)+  \varphi \left( Q \right)&=   \varphi \left( P \right)+  \varphi \left( Q\cap E_2 \right)+  \varphi \left( Q- E_2 \right)\\ 
          &=        \varphi \left( P\cup \left( Q-E_2 \right)  \right)+  \varphi \left( Q\cap E_2 \right)\\ 
           &=  \varphi \left( P\cup Q \right)    
        \end{aligned}
        \]
        \item  Let \(  S_{k} = \bigcup _{i= 1}^{k}E_{i}  \).  We proceed by finite induction and first note that the result is obviously true for \(  k= 1  \) . For \(  k> 1  \) assume that \(  S_{k}  \) is \(   \varphi   \)-measurable and that \[
         \varphi \left( A \right)\ge  \sum _{i = 1}^{k} \varphi \left( A\cap E_{i} \right)+  \varphi \left( A\cap S_{k}^{c} \right)   ,\quad \forall A\subseteq X
        \]   Then \[
        \begin{aligned}
         \varphi \left( A \right)&=  \varphi \left( A\cap E_{k+ 1} \right)+  \varphi \left( A\cap E_{k+ 1}^{c} \right)\\ 
          &=  \varphi \left( A\cap E_{k+ 1} \right)+  \varphi \left( A\cap E_{k+ 1}^{c}\cap S_{k}^{c} \right)+  \varphi \left( A\cap E_{k+ 1}^{c}\cap S_{k} \right)\\ 
           &=  \varphi \left( A\cap E_{k+ 1} \right)+  \varphi \left( A\cap S_{k+ 1}^{c} \right)+  \varphi \left( A\cap S_{k} \right)        
        \end{aligned}
        \]  Replace \(  A  \) by \(  A\cap S_{k}  \) , we have \[
         \varphi \left( A\cap S_{k} \right)\ge \sum _{i = 1}^{k} \varphi \left( A\cap S_{k}\cap E_{i} \right)+  \varphi \left( A\cap S_{k}\cap S_{k}^{c} \right)=     \sum _{i = 1}^{k} \varphi \left( A\cap E_{i} \right) 
        \]   Thus \[
         \varphi \left( A \right) \ge \sum _{i = 1}^{k+ 1} \varphi \left( A\cap E_{k+ 1} \right)+  \varphi \left( A\cap S_{k+ 1}^{c} \right)  
        \] Then the subadditivity gives that \[
         \varphi \left( A \right)\ge   \varphi \left( A\cap S_{k+ 1} \right)+  \varphi \left( A\cap S_{k+ 1}^{c} \right)   
        \] this , inturn , implies that \(  S_{k+ 1}  \) is \(   \varphi   \)-measurable. Since we now know  for every set \(  A\subseteq X  \) and for all positive integers \(  k  \) \[
         \varphi \left( A \right)\ge \sum _{i = 1}^{k} \varphi \left( A\cap E_{i} \right)+  \varphi \left( A\cap S_{k}^{c} \right)   
        \]     We have \[
         \varphi \left( A \right)\ge \sum _{i = 1}^{\infty} \varphi \left( A\cap E_{i} \right)+  \varphi \left( A\cap S^{c} \right)  \ge  \varphi \left( A\cap S \right)+  \varphi \left( A\cap S^{c} \right)   
        \] where \(  S= \bigcup _{i = 1}^{\infty}E_{k}  \) .  Thus \(  S  \) is measurable.   The subadditivity gives that \[
         \varphi \left( A \right)\le \sum _{i= 1}^{\infty} \varphi \left( A\cap E_{i} \right)+  \varphi \left( A\cap S^{c} \right)   
        \]Finally, we have \[
         \varphi \left( A \right)= \sum _{i= 1}^{\infty} \varphi \left( A\cap E_{i} \right)+  \varphi \left( A\cap S^{c} \right)   
        \] Replace \(  A  \) by \(  S  \) , we can get \[
         \varphi \left( S \right)= \sum _{i= 1}^{\infty} \varphi \left( E_{i} \right)  
        \]  
    \end{enumerate}
    
\end{proof}

\begin{lemma}{}{}
Let $\{E_i\}$ be a sequence of arbitrary sets. Then there exists a sequence of disjoint sets $\{A_i\}$ such that each $A_i \subset E_i$ and
\[
\bigcup_{i=1}^{\infty} E_i = \bigcup_{i=1}^{\infty} A_i.
\]
If each $E_i$ is $\varphi$-measurable, and so is $A_i$.
\end{lemma}
\begin{proof}
    Let \(  S_{k}= \bigcup _{i= 1}^{k}E_{i}  \) .  Let  \(  A_{k}= S_{k}\setminus S_{k-1} ,k \ge 1  \) , where \(  S_0   = \varnothing\) .  Then \[
    \bigcup _{i = 1}^{\infty}E_{i}= \bigcup _{i = 1}^{\infty}\left( S_{k}\setminus S_{k-1} \right)  = \bigcup _{i = 1}^{\infty}A_{i} .
    \]  If each \(  E_{i}  \) is \(   \varphi   \)-measurable.   \(  E_{i+ 1}\cup E_{i}= \left( E_{i+ 1}-E_{i} \right)\cup E_{i}    \) is measurable follows from the (ii), (iii) of the above theorem . Inductively , every \(  S_{k}=  \left( S_{k-1}\setminus E_{k} \right)\cup E_{k}   \)  is measurable.  Then  each \(  A_{k}   \) is measurable .  
\end{proof}

\begin{theorem}{}{7-24-2}
If $\{E_i\}$ is a sequence of \dfntxt{$\varphi$-measurable} sets in X, then $\cup_{i=1}^{\infty} E_i$ and $\cap_{i=1}^{\infty} E_i$ are \dfntxt{$\varphi$-measurable}.
\end{theorem}

\begin{proof}
    By the lemma above, there exists a sequence of disjoint \(   \varphi   \) - measurable sets \(  \left\{ A_{i} \right\}  \) such that each \(  A_{i}\subseteq E_{i}  \) , and \[
    \bigcup _{i = 1}^{\infty}E_{i}= \bigcup _{i = 1}^{\infty}A_{i}
    \]  It follows from Theorem \ref{thm:7-24-1} that \(  \bigcup _{i = 1}^{\infty}A_{i}  \) is measurable , then \(  \bigcup _{i = 1}^{\infty}E_{i}  \) is as well .   Since \(  X   \) is measurable , \(  E_{i}^{c}= X\setminus E_{i}  \) is measurable , then   \(  \bigcup _{i = 1}^{\infty}E_{i}^{c}  \) is measurable .  Thus \[
    \bigcap_{i= 1}^{\infty}E_{i}=  X\setminus \bigcup _{i= 1}^{\infty}E_{i}^{c} 
    \] is measurable . 
\end{proof}

\begin{definition}{}{}
A nonempty collection $\Sigma$ of sets $E$ satisfying the following two conditions is called a \dfntxt{$\sigma$-algebra}:
\begin{enumerate}
    \item[(i)] if $E \in \Sigma$, then $E^{c} \in \Sigma$;
    \item[(ii)] $\cup_{i=1}^{\infty} E_i \in \Sigma$ if each $E_i$ is in $\Sigma$.
\end{enumerate}
\end{definition}

\begin{definition}{}{}
In a topological space, the elements of the smallest $\sigma$-algebra that contains all open sets are called \dfntxt{Borel sets}. The term ``smallest'' is taken in the sense of inclusion.
\end{definition}
\begin{proof}
    We need to show that the ``smallest'' exists. 
\end{proof}

\begin{corollary}{}{}
If $\varphi$ is an \dfntxt{outer measure} on an arbitrary set $X$, then the class of $\varphi$-measurable sets forms a \dfntxt{$\sigma$-algebra}.
\end{corollary}
\begin{proof}
    Follows immediately from Theorem \ref{thm:7-24-1} and Theorem \ref{thm:7-24-2}. 
\end{proof}

\begin{corollary}{}{}
Suppose $\varphi$ is an outer measure on $X$ and $\{E_i\}$ a countable collection of $\varphi$-measurable sets.
\begin{enumerate}
    \item[(i)] If $E_1 \subset E_2$ with $\varphi(E_1) < \infty$, then
    \[
        \varphi(E_2 \setminus E_1) = \varphi(E_2) - \varphi(E_1).
    \]
    \item[(ii)] (Countable additivity) If $\{E_i\}$ is a disjoint sequence of sets, then
    \[
        \varphi(\bigcup_{i=1}^{\infty} E_i) = \sum_{i=1}^{\infty} \varphi(E_i).
    \]
    \item[(iii)] (Continuity from the left) If $\{E_i\}$ is an increasing sequence of sets, that is, if $E_i \subset E_{i+1}$ for each $i$, then
    \[
        \varphi(\bigcup_{i=1}^{\infty} E_i) = \varphi(\lim_{i \to \infty} E_i) = \lim_{i \to \infty} \varphi(E_i).
    \]
    \item[(iv)] (Continuity from the right) If $\{E_i\}$ is a decreasing sequence of sets, that is, if $E_i \supset E_{i+1}$ for each $i$, and if $\varphi(E_{i_0}) < \infty$ for some $i_0$, then
    \[
        \varphi(\bigcap_{i=1}^{\infty} E_i) = \varphi(\lim_{i \to \infty} E_i) = \lim_{i \to \infty} \varphi(E_i).
    \]
    \item[(v)] If $\{E_i\}$ is any sequence of $\varphi$-measurable sets, then
    \[
        \varphi(\liminf_{i \to \infty} E_i) \le \liminf_{i \to \infty} \varphi(E_i).
    \]
    \item[(vi)] If
    \[
        \varphi(\bigcup_{i=i_0}^{\infty} E_i) < \infty
    \]
    for some positive integer $i_0$, then
    \[
        \varphi(\limsup_{i \to \infty} E_i) \ge \limsup_{i \to \infty} \varphi(E_i).
    \]
    
\end{enumerate}
\end{corollary}

\begin{proof}
    (v)  Let \(  A_{j}=  \bigcap_{k= j}^{\infty}E_{k}   \) , then \[
    \liminf_{i\to \infty}E_{i}= \lim_{j\to \infty}A_{j}= \bigcup _{j = 1}^{\infty}A_{j}
    \] It follows from (iv) that \[
     \varphi \left( \liminf_{i\to \infty}E_{i} \right)= \lim_{j\to \infty} \varphi \left( A_{j} \right)  
    \] Since \(  A_{j}\subseteq  E_{j}  \), we have \(   \varphi \left( A_{j} \right)\le  \varphi \left( E_{j} \right)    \), then \[
    \lim_{j\to \infty} \varphi \left( A_{j} \right)\le \liminf_{i\to \infty} \varphi \left( E_{j} \right)  
    \]   Thus \[
     \varphi \left( \liminf_{i\to \infty}E_{i} \right)\le \liminf_{i\to \infty} \varphi \left( E_{i} \right)  
    \]
\end{proof}

\begin{definition}{}{}
    \begin{itemize}
        \item An outer measure \(   \varphi   \) on a topological space \(  X  \) is called a \dfntxt{Borel outer measure} if all Borel sets are \(   \varphi   \)-measurable. 
        \item A Borel outer measure is \dfntxt{finite} if \(   \varphi \left( X \right)    \) is finite.
        \item An outer measure \(   \varphi   \) on a set \(  X  \) is called \dfntxt{regular} if for each \(  A\subseteq X  \) there exists a \(   \varphi   \)-measurable set \(  B\supseteq A  \) such that \(   \varphi \left( B \right)=  \varphi \left( A \right)    \). 
        \item A \dfntxt{Borel rugular outer measure} is a Borel outer measure such that for each \(  A\subseteq X  \), there exists a Borel set \(  B  \) such that \(   \varphi \left( B \right)=  \varphi \left( A \right)    \).    
        \item A \dfntxt{Radon outer measure} is a Borel regular outer measure that is finite on compact sets.             
    \end{itemize}
    
\end{definition}
\section*{Exercises for Outer Measure}

\begin{exercise}{}{}
     If \(  X  \) is a metric space, fix \(   \varepsilon > 0  \). Let \(   \varphi \left( A \right)   \) be the smallest number of balls of radius \(   \varepsilon   \) that cover \(  A  \).      

Let \(   \varphi _{ \varepsilon }\left( A \right): =  \varphi \left( A \right)    \) denote the dependence on \(   \varepsilon   \), and define \[
\psi \left( A \right) : = \lim_{ \varepsilon \to 0 }  \varphi _{ \varepsilon }\left( A \right)  
\] What is \(  \psi \left( A \right)   \)  and what are the correspoding \(  \psi   \)-measurable sets. 
\end{exercise}
\begin{proof}
    If there are only finite points in \(  A  \). Let \[
    l : =  \min \left\{ d\left( x,y \right) :x,y \in A ,x \neq y\right\}
    \] Then for any \(  x,y \in A  \), \(  x \neq y  \). If \(  B_1,B_2  \) are ball of radius \( \frac{l }{2 }   \) such that \(  x\in B_1, y \in B_2  \).     Then \(  B_1\cap B_2= \varnothing  \).  But each point in \(  A  \) must covered by one of such ball, which shows that \(  \psi \left( A \right)\ge  \left| A \right|    \).    Since it is obvious that \(  \psi \left( A \right)\le \left| A \right|    \),  then \(  \psi \left( A \right)= \left| A \right|    \). 
    
    Now suppose that \(  A  \) is a infinite set. If \(  \psi \left( A \right)< \infty   \), say \(  \psi \left( A \right)= n   \) . Then for any sufficiently small \(   \varepsilon   \), there are balls \(  B_1,\cdots ,B_{n}  \)  with radius \(   \varepsilon   \) cover \(  A  \). Take any \(  n+ 1  \) different points \(  p_1,\cdots ,p_{n+ 1}  \)  of \(  A  \). Let  \(   \varepsilon = \frac{1}{2} \min \left\{ d\left( p_{i},p_{j} \right): i\neq j  \right\}  \), then \(   B_1,\cdots,B_n   \) cannot cover these \(  n+ 1  \) points, thus cannot cover  \(  A  \), which is a contradiction. Thus we have \(  \psi \left( A \right)= \infty   \) . 

    We find that \(  \psi   \) is just the counting measure on \(  X  \).Take any \(  E\subseteq X  \),and take any \(  A\subseteq X  \), if \(  A  \) is finite, then \(  A\cap E, A- E  \) are as well, and \(  \left| A \right|= \left| A\cap E \right|+ \left| A-E \right|     \)   , that is \(  \psi \left( A \right)= \psi \left( A\cap E \right)+ \psi \left( A-E \right)     \).  If \(  A  \) is infinite, then one of the \(  A\cap E  \)   and \(  A-E  \) is as well. Then \[
    \psi \left( A \right)= \infty= \psi \left( A\cap E \right)+ \psi \left( A-E \right)   
    \]The above shows that \(  E  \) is \(  \psi   \)-measurable. Thus all subset of \(  X  \) are \(  \psi   \)-measurable.        
\end{proof}

\begin{exercise}{}{}
    It was shown that \(   \varphi \left( E_2\setminus E_1 \right)=  \varphi  \left( E_2 \right)- \varphi \left( E_1 \right)     \) if \(  E_1\subseteq E_2  \) are \(   \varphi   \)-measurable with \(   \varphi \left( E_1 \right)< \infty   \). Prove that this result remains true if \(  E_2  \) is not assumed to be \(   \varphi   \)-measurable.      
\end{exercise}
\begin{proof}
   Since \(  E_1  \) is measurable,    \[
      \varphi \left( E_2 \right)=  \varphi \left( E_2\cap E_1 \right)+  \varphi \left( E_2-E_1 \right)   
     \] but \[
      \varphi \left( E_2\cap E_1 \right)=  \varphi \left( E_1 \right)  
     \]which completes the proof. 
\end{proof}

\section{Carath\'eodory Outer Measure}

\begin{definition}{Carath\'eodory outer measure}{}
    An outer measure \(   \varphi  \) defined on a metric space \(  \left( X,p \right)   \) is called a \dfntxt{Carath\'eodory outer measure} if \[
     \varphi \left( A\cup B \right)=  \varphi \left( A \right)+  \varphi \left( B \right)   
    \] whenever \(  A,B  \) are arbitrary subsets of \(  X  \) with \(  d\left( A,B \right)> 0   \). Where \[
    d\left( A,B \right): = \inf \left\{ \rho \left( a,b \right): a \in A, b\in B  \right\} 
    \]     
\end{definition}

\begin{theorem}{}{8-10-1}
    If \(   \varphi   \) is a Carath\'eodory outer measure on a metric space \(  X  \), then all closed sets are \(   \varphi   \)-measurable.    
\end{theorem}
\begin{proofsketch}
    用一列集合近似代替\(  C  \)是标准的想法. 证明的核心要素有: 1.\(  C  \)是闭集, 使得按距离向外逼近的行为是良好的. 2. 距离这个属性与Carath\'odory条件是相适应的.   此外, 将命题转化为极限 \[
    \lim_{i \to \infty}  \varphi \left( A\setminus C_{i} \right)=  \varphi \left( A\setminus C \right)  
    \]是标准的流程. 要证明该极限成立, 则需要一定的技巧. 这里将 \(  A\setminus C_{i}  \)和 \(  A\setminus C  \)相差的中间部分进行细分量化, 将原极限转化为 \[
\lim_{i \to \infty}    \sum _{j = i}^{\infty} \varphi \left( T_{j} \right) 
    \]的极限问题, 从而转化为级数的收敛性问题. 将集合隔一个分组, 变成两组两两相隔正距离的集合, 以便套用Carath\'eodory条件.   
\end{proofsketch}
\begin{proof}
    Let \(  C \subseteq X  \) be any closed sets. Take \(  A \subseteq X  \), only need to show that \[
     \varphi \left( A \right)\ge  \varphi \left( A\cap C \right) +  \varphi \left( A\setminus C \right)   
    \]We assume that \(   \varphi \left( A \right)< \infty   \).  Define \[
    C_{i}: = \left\{ x: d\left( x,C \right)\le \frac{1 }{i }   \right\}
    \] Since \(  d\left( A\setminus C_{i}, A\cap C \right)   > 0\),  \[
     \varphi \left( A \right)\ge   \varphi \left( \left( A\setminus C_{i} \right)\cup \left( A\cap C \right)   \right)=   \varphi \left( A\setminus C_{i} \right)+  \varphi \left( A\cap C \right)    
    \] Only need to show that \[
    \lim_{i \to \infty}  \varphi \left( A\setminus C_{i} \right)=  \varphi \left( A\setminus C \right)  
    \] Define \[
    T_{i}:=  \left\{ x \in A: \frac{1 }{i+ 1 }< d\left( x,C \right)\le \frac{1 }{i }    \right\}
    \], then \[
    A\setminus C\subseteq  \left( A\setminus C_{i} \right)\cup \bigcup _{j = i}^{\infty}T_{i} 
    \] Thus \[
     \varphi \left( A\setminus C \right)\le  \varphi \left( A\setminus C_{i} \right)+ \sum _{j = i}^{\infty} \varphi \left( T_{i} \right)   
    \] Once we have \(  \sum _{j= 1}^{\infty} \varphi \left( T_{i} \right)< \infty   \), then \(  \lim_{i\to \infty}\sum _{j = i}^{\infty} \varphi \left( T_{i} \right)= 0   \),  the conclusion follows from it. To show this, note that \(  \,\mathrm{d} \left( T_{i},T_{j} \right)> 0   \) whenever \(  \left| i-j \right|\ge 2   \). Then it by induction follows from Carath\'odory condition that \[
    \sum _{ j= 1}^{\infty} \varphi \left( T_{2i} \right) \le  \varphi \left( A \right) < \infty
    \] \[
    \sum _{j = 1}^{\infty} \varphi \left( T_{2j-1} \right)\le  \varphi \left( A \right) < \infty 
    \]    Thus \(  \sum _{j = 1}^{\infty} \varphi \left( T_{i} \right)< \infty   \), which completes the proof.  
\end{proof}

\begin{theorem}{}{8-10-2}
Suppose $\mathcal{F}$ is a family of subsets of a topological space X that contains all open and all closed subsets of X. Suppose also that $\mathcal{F}$ is closed under countable unions and countable intersections. Then $\mathcal{F}$ contains all Borel sets; that is, $\mathcal{B} \subset \mathcal{F}$.
\end{theorem}
\begin{note}
    可以将包含\(  \mathcal{B}  \)描述成对以下三种集合的封闭性: 1.开集 2. 补集 3.可数并. 

    如果我们事先没有对补集的封闭性, 可以通过加强1.3.的条件,使得当我们强制使得补集封闭后, 得到的新集族依旧满足1.3. 具体地, 我们希望原集族对于开集的补集,也就是闭集也封闭; 对于可数并的补集, 也就是可数交也封闭. 
\end{note}
\begin{proof}
    Let \[
    \mathcal{H}=  \mathcal{F}\cap \left\{ A:A^{c}\in \mathcal{F} \right\}
    \] Obseve that \(  \mathcal{H}  \) contains all closed sets. Moreover, it is easily seen that \(  \mathcal{H}  \) is closed under complementation and countable unions. Thus \(  \mathcal{H}  \) is a \(   \sigma   \)-algebra that contains the open sets and therefore contains all Borel sets.    
\end{proof}

\begin{theorem}{}{}
    If \(   \varphi   \) is a Carath\'eodory outer measure on a metric space \(  X  \), then the Borel sets of \(  X  \) are \(   \varphi   \)-measurable.    
\end{theorem}
\begin{proof}
    Since the collection of \(   \varphi   \)-measurable sets is itself a \(   \sigma   \)-algebra , it is closed under complements and countable union and countable intersection. Theorem \ref{thm:8-10-1} gives that \(   \varphi   \)-measurable sets contains closed sets, thus contains open sets as well.  Then  the conclusion follows from Theorem \ref{thm:8-10-2} .
\end{proof}

\section*{Exercises for this section}

\begin{exercise}{Smallest \(   \sigma   \)-algebra containing open sets }{}
    Prove that in every topological space $X$, there exists a smallest $\sigma$-algebra that contains all open sets in $X$. That is, prove that there is a $\sigma$-algebra $\Sigma$ that contains all open sets and has the property that if $\Sigma_1$ is another $\sigma$-algebra containing all open sets, then $\Sigma \subset \Sigma_1$. In particular, for $X = \mathbb{R}^n$, note that there is a smallest $\sigma$-algebra that contains all the closed sets in $\mathbb{R}^n$.
\end{exercise}

\begin{proof}
    Let \[
     \Sigma : = \bigcap_{\mathcal{M}_{\alpha } \text{ is a } \sigma \text{ algebra containing }X }\mathcal{M}_{\alpha } 
    \] Then it is obvious that \(   \Sigma   \) is closed under complement,  countable union, and contains all open sets. Thus \(   \Sigma   \) is a \(   \sigma   \)-algebra containing all open sets.   
\end{proof}

\begin{exercise}{}{}
     In a topological space $X$ the family of Borel sets, $\mathcal{B}$, is by definition the $\sigma$-algebra generated by the closed sets. The method below is another way of describing the Borel sets using transfinite induction. You are to fill in the necessary steps:
    \begin{enumerate}
        \item For an arbitrary family $\mathcal{F}$ of sets, let
        \[
            \mathcal{F}^* = \left\{ \bigcup_{k=1}^{\infty} E_k : \text{ where either } E_i \in \mathcal{F} \text{ or } \tilde{E}_i \in \mathcal{F} \text{ for all } i \in \mathbb{N} \right\}.
        \]
        Let $\Omega$ denote the smallest uncountable ordinal. We will use transfinite induction to define a family $\mathcal{E}_\alpha$ for each $\alpha < \Omega$.
        \item Let $\mathcal{E}_0 :=$ all closed sets $:= \mathcal{K}$. Now choose $\alpha < \Omega$ and assume that $\mathcal{E}_\beta$ has been defined for each $\beta$ such that $0 \le \beta < \alpha$. Define
        \[
            \mathcal{E}_\alpha := \left( \bigcup_{0 \le \beta < \alpha} \mathcal{E}_\beta \right)^*
        \]
        and define
        \[
            \mathcal{A} := \bigcup_{0 \le \alpha < \Omega} \mathcal{E}_\alpha.
        \]
        \item Show that each $\mathcal{E}_\alpha \subseteq \mathcal{B}$.
        \item Show that $\mathcal{A} \subset \mathcal{B}$.
        \item Now show that $\mathcal{A}$ is a $\sigma$-algebra to conclude that $\mathcal{A} = \mathcal{B}$.
        \begin{enumerate}
            \item Show that $\emptyset, X \in \mathcal{A}$.
            \item Let $A \in \mathcal{A} \implies A \in \mathcal{E}_\alpha$ for some $\alpha < \Omega$. Show that this implies $\tilde{A} \in \mathcal{E}_\alpha^* \subset \mathcal{E}_\beta$ for every $\beta > \alpha$.
            \item Conclude that $\tilde{A} \in \mathcal{A}$ and thus conclude that $\mathcal{A}$ is closed under complementation.
            \item Now let $\{A_k\}$ be a sequence in $\mathcal{A}$. Show that $\bigcup_{k=1}^{\infty} A_k \in \mathcal{A}$.
            (Hint: Each $A_k$ is in $\mathcal{A}_{\alpha_k}$ for some $\alpha_k < \Omega$. We know that there is $\beta < \Omega$ such that $\beta > \alpha_k$ for each $k \in \mathbb{N}$.)
        \end{enumerate}
    \end{enumerate}
\end{exercise}

\begin{proof}
    3. Since \(  \mathcal{E}_{0} \subseteq  \mathcal{B}  \) , only need to show that if \(  \mathcal{E}_{\beta } \subseteq \mathcal{B}  \) for all \(  0\le \beta < \alpha   \) , then \(  \mathcal{E}_{\alpha }\subseteq \mathcal{B}  \). Take any \(  E \in \mathcal{E}_{\alpha }  \) , then there are \(   E_1,\cdots,E_k ,\cdots   \), such that \(  E_{i} \in \bigcup _{0 \le \beta < \alpha }\mathcal{E}_{\beta }  \), or \(  E_{i}^{c}\in \bigcup _{0 \le \beta < \alpha } \mathcal{E}_{\beta }  \) , and \(  E=  \bigcup _{k= 1}^{\infty}E_{k}  \).  We suppose that \(  E_{i} \in \mathcal{E}_{\beta _{i}}  \) or \(  E_{i}^{c} \in \mathcal{E}_{\beta _{i}}  \). Since \(  \mathcal{E}_{\beta _{i}} \subseteq \mathcal{B}  \) for all \(  \beta _{i}  \), and since \(  \mathcal{B}  \) is closed under complement, then \(  E_{i} \in \mathcal{B}  \). Since \(  \mathcal{B}  \) is closed under countable union, \(  E= \bigcup _{k = 1}^{\infty}E_{k} \in \mathcal{B}  \). Thus \(  \mathcal{E}_{\alpha } \subseteq \mathcal{B}  \). 

    4. We have shown that \(  \mathcal{E}_{\alpha } \subseteq \mathcal{B}  \) for all \(  \alpha <  \Omega   \),  then it is immediate that \(  A\subseteq \mathcal{B}  \).  

    5. (a)  Since \(  \varnothing \in \mathcal{E}_{0}  \) and \(  X= \left( \varnothing \right)^{c}\in \mathcal{E}_{1}= \left( \mathcal{E}_{0} \right)^{*}   \)  , then \(  \varnothing,X \in \mathcal{A}  \).

    (b)  Write \[
    \mathcal{E}_{\alpha }^{*}= \left\{ \bigcup _{k= 1}^{\infty}E_{k}: \text{  where either }E_{i} \in \mathcal{E}_{\alpha } \text{ or } E_{i}^{c}\in \mathcal{E}_{\alpha }\text{ for all } i\in \mathbb{N}  \right\}
    \]We take \(  E_1 = A^{c} \), then \(  E_1^{c}= A\in \mathcal{E}_{\alpha }  \), and take \(  E_{2}= \cdots = \varnothing  \).Then \(  A^{c}= \bigcup _{k = 1}^{\infty}E_{k} \in \mathcal{E}_{\alpha }^{*}  \). If \(  \beta  >  \alpha   \), then  \[
   \mathcal{E}_{\alpha }\subseteq  \bigcup _{0 \le  \gamma < \beta }\mathcal{E}_{ \gamma  }
    \]   Let \(  *  \) act on both side, then it follows from the definition of \(  *  \) that \[
    \left( \mathcal{E}_{\alpha } \right)^{*}\subseteq \left( \bigcup _{0\le  \gamma < \beta }\mathcal{E}_{ \gamma } \right)^{*}= \mathcal{E}_{\beta }  
    \]  Thus \(  A^{c} \in \mathcal{E}_{\alpha }^{*}\subseteq \mathcal{E}_{\beta }  \) for every \(  \beta > \alpha   \). 
    
    (c) It is immediate from (b) that \(  A^{c} \in \mathcal{E}_{\beta }\subseteq \mathcal{A}  \) .

    (d) Suppose that \(  A_{k} \in \mathcal{E}_{\alpha _{k}}  \) for some \(  \alpha _{k}<  \Omega   \). Let \[
     \gamma : = \bigcup _{k \in \mathbb{N} }\alpha _{k},\quad \beta : =  \gamma + 1
    \] Then \(   \gamma ,\beta <  \Omega   \), and \(  \alpha _1 ,\cdots ,\alpha _{k}< \beta   \). Then \[
    \bigcup _{k = 1}^{\infty}A_{k}\in \left( \bigcup _{k = 1}^{\infty} \mathcal{E}_{\alpha _{k}} \right)^{*}\subseteq \mathcal{E}_{\beta }\subseteq \mathcal{A} 
    \]   
    
\end{proof}

\begin{exercise}*{}{}
      An outer measure \(   \varphi   \) on a space \(  X  \) is called \(   \sigma   \)-finite if there exists a countable number of sets \(  A_{i}  \) with \(   \varphi \left( A_{i} \right)< \infty   \) such that \(  X\subseteq \bigcup _{i= 1}^{\infty}A_{i}  \). Assuming that \(   \varphi   \) is a \(   \sigma   \)-finite Borel regular outer measure on a metric space \(  X  \), prove that \(  E\subseteq X  \) is \(   \varphi   \)-measurable if and only if there exists an \(  F_{ \sigma }   \) set \(  F\subseteq E  \) such that \(   \varphi \left( E\setminus F \right)= 0   \).       
\end{exercise}
\begin{proof}
    Let \(  E_{i}: = E\cap A_{i}  \). Then \[
    E=  \bigcup _{i = 1}^{\infty}E_{i},\quad  \varphi \left( E_{i} \right)\le  \varphi \left( A_{i} \right)< \infty  
    \] Once we show that there is a closed set \(  F_{i}\subseteq E_{i}  \), such that \(   \varphi \left( E_{i}\setminus F_{i} \right)= 0   \) for all \(  i  \), then let \(  F= \bigcup _{i = 1}^{\infty}F_{i}\in F_{ \sigma }  \), it follows that  \[
     \varphi \left( E\setminus F \right)\le  \varphi \left( \bigcup \left( E\setminus F_{i} \right)  \right)\le  \varphi \left( \bigcup \left( E_{i}\setminus F_{i} \right)  \right)\le \sum _{i} \varphi \left( E_{i}\setminus F_{i} \right)= 0    
    \]   
\end{proof}


\begin{exercise}{}{}
Let $\varphi$ be an outer measure on a space $X$. Suppose $A \subset X$ is an arbitrary set with $\varphi(A) < \infty$ such that there exists a $\varphi$-measurable set $E \supset A$ with $\varphi(E) = \varphi(A)$. Prove that $\varphi(A \cap B) = \varphi(E \cap B)$ for every $\varphi$-measurable set $B$.
\end{exercise}
\begin{proof}
   \[
      \varphi \left( A \right)=  \varphi \left( E \right)=  \varphi \left( E\cap B \right)+  \varphi \left( E\setminus B \right)   
     \]  \[
    \begin{aligned}
      \varphi \left( A\cap B \right)\ge  \varphi \left( A \right)- \varphi \left( A\setminus B \right)&=  \varphi \left( E \right)- \varphi \left( A\setminus B \right)  \\ 
       &=  \varphi \left( E\cap B \right)+  \varphi \left( E\setminus B \right)- \varphi \left( A\setminus B \right)\\ 
        &\ge  \varphi \left( E\cap B \right)    
    \end{aligned}    
     \] Since \(  A\cap B\subseteq E\cap B  \), \(   \varphi \left( A\cap B \right)\le  \varphi \left( E\cap B \right)    \). Finally, we have \(   \varphi \left( A\cap B \right)=  \varphi \left( E\cap B \right)    \)   
\end{proof}


\begin{exercise}{}{}
In $\mathbb{R}^2$, find two disjoint closed sets $A$ and $B$ such that $\mathbf{d}(A, B) = 0$. Show that this is not possible if one of the sets is compact.
\end{exercise}

\begin{proof}
    \begin{enumerate}
        \item Let \(  A  \) be the graph of \(  y= \frac{1 }{x }   \). Let \(  B = \left\{ \left( x,y \right): y= 0  \right\}  \). Then \(  \left\{ \left( n,\frac{1 }{n }  \right): n\in \mathbb{N}   \right\}  \subseteq A\) , \(  \left\{ \left( n,0 \right): n \in \mathbb{N}   \right\}\subseteq B  \), then \[
        \mathbf{d}\left( A,B \right)\le \inf \left\{ d\left( \left( n,\frac{1 }{n }  \right), \left( n,0 \right)   \right)  : n \in \mathbb{N} \right\} = \lim_{n\to \infty}\frac{1 }{n }= 0  
        \]
        \item If \(  A  \) is compact, \(  B  \) is closed. Since \(  \mathbf{d}\left( A,B \right)= 0   \), there are convergent sequences \(  \left\{ a_{n} \right\}\subseteq A  \) and \(  \left\{ b_{n} \right\}\subseteq B  \) such that \(  \lim_{n\to \infty}d\left( a_{n},b_{n} \right)= 0   \).   Since \(  A  \) is bounded and closed, there is a  cluster point \(   a \in A  \) of \(  \left\{ a_{n} \right\}  \). Suppose that \(  \lim_{k\to \infty}a_{n_{k}}= a  \), then \(  \lim_{k\to \infty}b_{n_{k}}= a\in B  \). \(  A\cap B\neq \varnothing  \), which is a contradiction.      
    \end{enumerate}
    
\end{proof}
\begin{exercise}{}{}
Let $\varphi$ be an outer Carathéodory measure on $\mathbb{R}$ and let $f(x) := \varphi(I_x)$, where $I_x$ is an open interval of fixed length centered at $x$. Prove that $f$ is lower semicontinuous. What can you say about $f$ if $I_x$ is taken as a closed interval? Prove the analogous result in $\mathbb{R}^n$; that is, let $f(x) := \varphi(B(x, a))$, where $B(x, a)$ is the open ball with fixed radius $a$ centered at $x$.
\end{exercise}
\begin{exercise}{}{}
In a metric space $X$, prove that $\operatorname{dist}(A, B) = d(\bar{A}, \bar{B})$ for arbitrary sets $A, B \in X$.
\end{exercise}
\begin{exercise}{}{}
Let $A$ be a non-Borel subset of $\mathbb{R}^n$ and define for each subset $E$,
\[\begin{aligned}
\varphi(E) = \begin{cases} 0 & \text{if } E \subset A, \\ \infty & \text{if } E \setminus A \neq \emptyset. \end{cases}
\end{aligned}\]
Prove that $\varphi$ is an outer measure that is not Borel regular.
\end{exercise}
\begin{exercise}{}{}
Let $\mathcal{M}$ denote the class of $\varphi$-measurable sets of an outer measure $\varphi$ defined on a set $X$. If $\varphi(X) < \infty$, prove that the family
\[
\mathcal{F} := \{A \in \mathcal{M} : \varphi(A) > 0\}
\]
is at most countable.
\end{exercise}


\section{Lebesgue Measure}

\begin{theorem}{}{}
Suppose each edge $I_k = [a_k, b_k]$ of an $n$-dimensional interval $I$ is partitioned into $\alpha_k$ subintervals. The products of these intervals produce a partition of $I$ into $\beta := \alpha_1 \cdot \alpha_2 \cdots \alpha_n$ subintervals $I_i$ and
\[
v(I) = \sum_{i=1}^{\beta} v(I_i).
\]
\end{theorem}


\begin{theorem}{}{}
For each interval $I$ and each $\epsilon > 0$, there exists an interval $J$ whose interior contains $I$ and
\[
v(J) < v(I) + \epsilon.
\]
\end{theorem}


\begin{definition}{Lebesgue outer measure}{}
    The \dfntxt{Lebesgue outer measure} of an arbitrary set \(  E\subseteq \mathbb{R} ^{n}  \), denoted by \(   \lambda ^{*}\left( E \right)   \), is defined by \[
     \lambda ^{*}\left( E \right): = \inf \left\{ \sum _{k= 1}^{\infty}v\left( I _{k} \right)  \right\}, 
    \] where the infimum is taken over all countable collections of closed intervals \(  I _{k}  \) such that \[
    E\subseteq \bigcup _{k = 1}^{\infty}I _{k}. 
    \]  
\end{definition}

\begin{remark}
   \begin{itemize}
    \item  It may be necessary at times to emphasize the dimension of the Euclidean space in which Lebesgue outer measure if defined. When clarification is needed, we will write \(   \lambda _{n}^{*}\left( E \right)   \) in plce place of \(   \lambda ^{*}\left( E \right)   \).   
    \item We have not prove that \(  \dfntxt{Lebesgue outer measure}  \) is actually a outer measure till now. 
   \end{itemize}
   
\end{remark}


\begin{theorem}{}{}
    For a closed interval \(  I\subseteq \mathbb{R} ^{n}  \), \(   \lambda ^{*}\left( I \right)= v\left( I \right)    \).   
\end{theorem}

\begin{proof}
    The inequality \(   \lambda ^{*}\left( I \right)\le v\left( I \right)    \) holds, since \(  I  \) itself givs a element of \(  \left\{ \sum _{k = 1}^{\infty}v\left( I _{k} \right)  \right\}  \)   in the definition of \(   \lambda ^{*}\left( I \right)   \). 

    To prove the opposite inequality, choose \(   \varepsilon > 0  \), and a sequence of closed intervals \(  I_1,I_2,\cdots   \) such that \[
    I\subseteq \bigcup _{k = 1}^{\infty}I _{k},\quad \sum _{ k= 1}^{\infty}v\left( I _{k} \right) \le  \lambda ^{*}\left( E \right)+  \varepsilon   
    \]  We can choose closed intervals \(  J_{k}  \) such that \[
    v\left( J_{k} \right)\le v\left( I _{k} \right)+ \frac{ \varepsilon  }{2^{k} },\quad  I _{k}\subseteq J_{k}^{\circ}   
    \]Let \(  \mathcal{F}  : =  \left\{ J_{k}^{\circ}:k \in \mathbb{N}  \right\}\).    Then \(  \mathcal{F}  \) is an open cover of \(  I  \), since \(  I  \) is compact, it has a lebesgue number \(  \eta   \).    There is a finitely many partition \(   K_1,\cdots,K_m   \) of \(  I  \) such that each  with diameter less than \(  \eta   \). Then \(  K_{i}  \) is contained in some  element of \(  \mathcal{F}  \), we say \(  J_{k_{i}}  \)    althoung two \(  K_{i}  \) may be contained in a same \(  J_{k_{i}}  \).  Let \(  N_{m}  \) be the smallest number of \(  J_{k_{i}}  \) containing \(  K_{i}  \), then \[
    v\left( I \right)= \sum _{i= 1}^{m}v\left( K_{i} \right)\le \sum _{i = 1}^{N_{m}}v\left( J_{k_{i}} \right)\le \sum _{ k = 1}^{\infty}v\left( J_{k} \right)    \le   \lambda ^{*}\left( E \right)+ 2 \varepsilon  
    \]     Since \(   \varepsilon   \) is arbitrary, \(  v\left( I \right)\le  \lambda ^{*}\left( E \right)    \).   
\end{proof}

\begin{theorem}{}{}
    Lebesgue outermeasure, \(   \lambda ^{*}  \), defined on \(  \mathbb{R} ^{n}  \) is a Carath\'eorody outer measure.  
\end{theorem}

\begin{proof}
    We first verify that \(   \lambda ^{*}  \) is an outer measure. We only show the subadditivity of \(   \lambda ^{*}  \), the other conditions are obvious. Let \( \left\{ A_{i} \right\}  \) be a countable collection of arbitrary subsets of \(  \mathbb{R} ^{n}  \). Denote \(  A: = \bigcup _{i = 1}^{\infty}A_{i}  \).       There is a countable family of closed intervals \(  \left\{ I _{j}^{\left( i \right) } \right\}_{j \in \mathbb{N} }  \) such that \[
    A_{i}\subseteq \bigcup _{j = 1}^{\infty} I^{\left( i \right) }_{j},\quad   \sum _{j = 1}^{\infty}v\left( I _{j}^{\left( i \right) } \right)<  \lambda ^{*}\left( A_{i} \right)+ \frac{ \varepsilon  }{2^{i} }   
    \] Collection of closed interval \(  \left\{ I _{j}^{\left( i \right) } \right\}_{i,j\in \mathbb{N} }  \)   can be aranged to a sequence of closed invervals, such that \[
    A\subseteq  \bigcup _{i,j\in \mathbb{N} } I^{\left( i \right) }_{j}
    \] Thus \[
     \lambda ^{*}\left( A \right)\le \sum _{i,j}v\left( I _{j}^{\left( i \right) } \right)\le \sum _{i = 1}^{\infty} \lambda ^{*}\left( A_{i} \right)+  \varepsilon    
    \] Since \(   \varepsilon > 0  \) is arbitrary, the countable subadditivity of \(   \lambda ^{*}  \) is established.      

    To show that \(   \lambda ^{*}  \) is Carath\'eorody outer measure. We take any \(   \varepsilon > 0  \) , and take a countable conllection of closed intervals \(  I_1,I_2,\cdots   \)   such that \[
    A\cup B\subseteq \bigcup _{ k= 1}^{\infty}I _{k},\quad  \sum _{k = 1}^{\infty}v\left( I _{k} \right)\subseteq  \lambda ^{*}\left( A\cup B \right).   
    \] We can subdividing each \(  I _{k}  \) into smaller invervals if necessary, we may assume that each interval \(   I _{k}  \) which diameter less that \(  \mathbf{d}\left( A,B \right)   \). Then \(  \left\{  I _{k} \right\}  \) can be seperated into two subfamiles \(  \left\{ I _{k}^{\prime}  \right\}  \) and \(  \left\{ I _{k}^{\prime \prime}  \right\}  \). The members in the former have empty intersection with \(  B  \) , and that  in the later have empty intersection with \(  A  \). \(  \left\{ I _{k} \right\}= \left\{ I _{k}^{\prime}  \right\} \sqcup \left\{ I _{k}^{\prime \prime}  \right\}  \). Thus \[
    A\subseteq \bigcup _{k = 1}^{\infty}\left\{ I _{k}^{\prime}  \right\},\quad B\subseteq \bigcup _{k= 1}^{\infty} \left\{ I _{k}^{\prime \prime}  \right\}
    \]          We have \[
    \begin{aligned}
     \lambda ^{*}\left( A \right)+  \lambda ^{*}\left( B \right)&\le  \sum _{k= 1}^{\infty}v\left( I _{k}^{\prime}  \right)+ \sum _{k = 1}^{\infty}v\left( I _{k}^{\prime \prime}  \right)     \\ 
      & = \sum _{ k= 1}^{\infty}v\left( I _{k} \right)\\ 
       &<   \lambda ^{*}\left( A\cup B \right)+  \varepsilon   
    \end{aligned}
    \] Since \(   \varepsilon > 0  \) is arbitrary, we have \[
     \lambda ^{*}\left( A \right)+  \lambda ^{*}\left( B \right)\le  \lambda ^{*}\left( A\cup B \right)   
    \] 
\end{proof}

\begin{corollary}{}{}
    All Borel sets in \(  \mathbb{R} ^{n}  \) are Lebesgue measurable. In particular, each open set and each closed set are Lebesgue measurable.  
\end{corollary}

\begin{definition}{Lebesgue Measure}{}
    Donote by \(   \lambda   \)  the set function obtained by restricting \(   \lambda ^{*}  \) to the family of Lebesgue measurable sets. \(   \lambda   \) is called \dfntxt{Lebesgue measure}.   
\end{definition}

\begin{corollary}{}{}
    Let \(  J  \) be an opoen interval in \(  \mathbb{R} ^{n}  \). Then \[
     \lambda \left( J \right)= \mathrm{vol}\left( J \right)  
    \]  
\end{corollary}
\begin{proof}
    There is a sequence of closed intervals \(  I _{k}  \)  such taht \(  J= \bigcup _{ k= 1}^{\infty}I _{k}  \), then \[
     \lambda \left( J \right)= \lim_{k\to \infty} \lambda \left( I _{k} \right)= \lim_{k\to \infty} \mathrm{vol}\left( I _{k} \right)= \mathrm{vol}\left( J \right)    
    \] 
\end{proof}
\end{document}